\documentclass{article}
\usepackage[utf8]{inputenc}
\usepackage{amsmath}
\usepackage{geometry}
\geometry{margin=2cm}
\begin{document}

\section*{Análise da Questão 153 — ENEM 2024 — Matemática}

\subsection*{Enunciado}
Um instituto de pesquisa constatou que, nos últimos dez anos, o crescimento populacional de uma cidade foi de 135,25\%. Qual é a representação decimal da taxa percentual desse crescimento populacional?

\section{Solução Técnica}

A taxa percentual de crescimento populacional é dada como 135,25\%.

\subsection{Estratégia}
Para converter uma taxa percentual para sua representação decimal, utilizamos a relação:

\begin{equation}
\text{taxa decimal} = \frac{\text{taxa percentual}}{100}.
\end{equation}

\subsection{Passos}
\begin{enumerate}
  \item Identificar que a taxa percentual é 135,25\%.
  \item Converter o percentual para decimal dividindo por 100.
  \item Calcular:
  \begin{equation}
  135,25 \div 100 = 1,3525.
  \end{equation}
  \item Confirmar que 1,3525 é a representação decimal correta.
\end{enumerate}

\noindent Portanto, a representação decimal correta do crescimento populacional é \textbf{1,3525}.

\noindent \textbf{Nota:} Esta é a alternativa \textbf{D} na questão original.

\section{Explicação Didática}

\subsection{Explicação resumida}
Para converter uma taxa percentual em sua forma decimal, basta dividir o valor percentual por 100. Assim, 135,25\% corresponde a 1,3525 em forma decimal.

\subsection{Passo a passo}
\begin{enumerate}
  \item Identifique o valor percentual fornecido: 135,25\%.
  \item Lembre-se que porcentagem significa "por 100", portanto, para converter para decimal, divida o valor por 100.
  \item Calcule: 135,25 \div 100 = 1,3525.
  \item Assim, a representação decimal do crescimento populacional é 1,3525.
\end{enumerate}

\subsection{Erros comuns}
\begin{itemize}
  \item Confundir valor percentual com valor decimal, usando diretamente 135,25 como resposta.
  \item Multiplicar por 100 em vez de dividir, obtendo 13.525.
  \item Dividir por 1000 ou outra potência incorreta de 10, gerando 0,13525.
  \item Deslocar a vírgula incorretamente, por exemplo, obtendo 13,525 ao invés de 1,3525.
\end{itemize}

\subsection{Dicas para o ensino}
\begin{itemize}
  \item Explique o conceito de porcentagem como uma fração por 100 para enfatizar a divisão.
  \item Use exemplos simples, como 50\% = 0,5, para facilitar a compreensão antes de avançar para números maiores.
  \item Incentive os alunos a praticar o deslocamento da vírgula nos números.
  \item Mostre os erros comuns para que eles possam reconhecê-los e evitá-los.
\end{itemize}

\section{Análise dos Distratores}

\begin{itemize}

  \item \textbf{A) 13 525,0}
    \begin{itemize}
      \item \textit{Possível raciocínio:} Erro ao multiplicar por 100 em vez de dividir.
      \item \textit{Tipo de erro:} Escala decimal incorreta.
      \item \textit{Por que parece plausível:} O valor tem a mesma magnitude do percentual original, mas sem a conversão necessária.
      \item \textit{Por que está errado:} 135,25\% corresponde a 1,3525 decimal, não 13.525.
    \end{itemize}

  \item \textbf{B) 135,25}
    \begin{itemize}
      \item \textit{Possível raciocínio:} Confusão entre valor percentual e decimal, usando diretamente o dado do enunciado.
      \item \textit{Tipo de erro:} Confusão entre porcentagem e número decimal.
      \item \textit{Por que parece plausível:} Apresenta o valor dado na questão sem modificação.
      \item \textit{Por que está errado:} A conversão para decimal exige a divisão por 100.
    \end{itemize}

  \item \textbf{C) 13,525}
    \begin{itemize}
      \item \textit{Possível raciocínio:} Deslocamento incorreto da vírgula, apenas uma casa para a esquerda.
      \item \textit{Tipo de erro:} Erro no posicionamento decimal.
      \item \textit{Por que parece plausível:} Parece ser um resultado de conversão parcial.
      \item \textit{Por que está errado:} O deslocamento correto é duas casas à esquerda, gerando 1,3525.
    \end{itemize}

  \item \textbf{E) 0,13525}
    \begin{itemize}
      \item \textit{Possível raciocínio:} Divisão por 1000 ao invés de 100.
      \item \textit{Tipo de erro:} Divisão por potência errada de 10.
      \item \textit{Por que parece plausível:} Resultado similar a uma conversão decimal, porém excessiva.
      \item \textit{Por que está errado:} A conversão correta pede divisão por 100.
    \end{itemize}

\end{itemize}

\section{Perfil do Item}

\subsection{Habilidades Sobretendidas}
\begin{itemize}
  \item Habilidade principal: Interpretação de porcentagem e conversão para número decimal.
  \item Habilidades secundárias:
    \begin{itemize}
      \item Operações com números decimais.
      \item Compreensão das diferentes representações numéricas (porcentual e decimal).
    \end{itemize}
\end{itemize}

\subsection{Demanda Cognitiva}
Nível baixo a médio.

\subsection{Dificuldade Estimada}
Baixa.

\subsection{Exigências da Questão}
\begin{itemize}
  \item Não necessita de interpretação visual complexa.
  \item Não requer modelagem.
  \item Requer cálculo simples e compreensão do conceito de porcentagem.
\end{itemize}

\subsection{Observações}
\begin{itemize}
  \item Trata-se de questão direta sobre conversão de valores percentuais para decimais.
  \item Explora conhecimento fundamental de porcentagem e conversão numérica.
\end{itemize}

\end{document}